% Immutable cited-input supplement. Target proofs remain embargoed.
% Target: wbII:test:cdp-comparison; tag 03IV.
% Frozen source SHA256 74497bfe327fae25c9cdcb17334b3bf42a9793a9df7c0a4404d97568f2f760f0.
% CompendiumRef is resolved against the cited S6_09 source edition, not silently identified with the current series statement.

% Exact cited input: comp:I / wb:thm:candidate-global-geometry
\begin{theorem}[global analytic construction]
\label{wb:thm:candidate-global-geometry}
\ Gluing the period family over \(B^\circ\) to
\(N_0,N_1,N_2\) by the preceding punctured-disc biholomorphisms produces a
connected compact Hausdorff complex threefold \(X\) and a proper surjective
holomorphic map
\[
 f:X\longrightarrow B\simeq\mathbb P^1
\]
with connected fibres.  The fibre over \(p_0\) is exactly the reduced
nonnormal surface \(W\) of \cref{wb:thm:candidate-cusp-geometry}; those over
\(p_1,p_2\) have scheme multiplicities three and four.  The open embedding
\(N_0\hookrightarrow X\) identifies canonical bundles, so
\begin{equation}
 \omega_X|_{N_0}\simeq\omega_{N_0}\simeq\mathcal O_{N_0}.
\label{wb:eq:candidate-global-canonical-on-cusp}
\end{equation}
\end{theorem}

% Exact cited input: comp:I / wb:thm:candidate-cusp-geometry
\begin{theorem}[the cusp quotient and its central fibre]
\label{wb:thm:candidate-cusp-geometry}
\ For the period functions and holomorphic cusp matrix \(C(t)\)
constructed above, and sufficiently small \(\epsilon>0\), the maps
\[
 \Psi_{\bar\lambda}(x,t)=
 \bigl(c_{\bar\lambda}(t)t^{B_0\bar\lambda}x,t\bigr),
 \qquad
 c_{\bar\lambda}(t)=\exp(2\pi iC(t)\bar\lambda),
\tag{2.5}
\]
extend to a free properly discontinuous action of
\(\bar\Lambda\simeq\mathbb Z^2\) on
\(Y_\epsilon=\{|t|<\epsilon\}\).  The quotient
\[
 f_0:N_0=Y_\epsilon/\bar\Lambda\longrightarrow\Delta_\epsilon
\]
is a proper holomorphic map from a connected smooth complex threefold.
Its nonzero fibres are the tori of \cref{wb:prop:period-input-audit}.  Its full
scheme-theoretic central fibre \(W\) is reduced and has local equations
\[
 z_0,\qquad z_0z_1,\qquad z_0z_1z_2.
\tag{2.6}
\]
The surface \(W\) is irreducible and nonnormal.  Its normalization is
\(\nu:E_0\simeq\mathrm{dP}_6\to W\), obtained by identifying opposite sides
of the toric hexagon; its double locus consists of three rational curves
through exactly two triple points.  Thus \(W\) is a normal-crossings surface,
but it is not a simple-normal-crossings divisor.  Moreover
\[
 \omega_{N_0}\simeq\mathcal O_{N_0}.
\tag{2.6a}
\]
\end{theorem}

% Exact cited input: comp:I / wb:prop:period-input-audit
\begin{proposition}[existence and non-degeneracy of the period matrix]
\label{wb:prop:period-input-audit}
\ There are holomorphic functions
\[
 \tau:\mathfrak h_z\to\mathfrak h,
 \qquad \mu,\beta:\mathfrak h_z\to\mathbb C
\]
satisfying
\[
\begin{array}{lll}
\tau\circ g_0=\tau-1,&\mu\circ g_0=\mu,&\beta\circ g_0=\beta+1,\\
\tau\circ g_2=-1/\tau,&\mu\circ g_2=1+\mu/\tau,
 &\beta\circ g_2=\beta-3-6\mu^2/\tau,
\end{array}
\tag{2.1}
\]
and the corresponding order-three transformation laws obtained from
\(\tau\circ g_1=(\tau-1)/\tau\).  They may be chosen so that
\[
 \operatorname{Im}\tau>0,
 \qquad
 D:=\operatorname{Im}\beta-
 \frac{6(\operatorname{Im}\mu)^2}{\operatorname{Im}\tau}<0.
\tag{2.2}
\]
Consequently
\[
 \Pi=[Z\mid I_2],\qquad
 Z=\begin{pmatrix}6\mu&\tau\\ \beta&\mu\end{pmatrix}
\tag{2.3}
\]
maps the marked rank-four lattice \(\Lambda\) isomorphically onto a lattice
in \(\mathbb C^2\) at every point.  The quotients
\(\mathbb C^2/\Pi(z)\Lambda\) form a holomorphic family of compact
two-dimensional complex tori over \(\mathfrak h_z\), equivariant for
\(\Delta\).
\end{proposition}

